\section{Discussion}
\label{sec:discussion}

\subsection{Assembly Identity as a TMRCA Estimator}

The central finding is that sequence identity from pangenome alignments provides a practical, theoretically grounded estimator of pairwise TMRCA. As established in Section~\ref{sec:introduction}, four independent research programs converge on the superiority of TMRCA-based relatedness. Our per-window identity matrix constitutes a local eGRM estimated directly from assemblies, bypassing both VCF generation and ARG inference. Formally, the identity matrix is an affine transformation of the Branch GRM: $G_w = (1-2\mu\bar{H}_w)\mathbf{11}' + 2\mu\cdot\text{bGRM}_w + E_w$, with spectral equivalence (Davis-Kahan angular deviation $<$1\textdegree) ensuring that identity-based PCA reproduces Branch GRM structure. Per-window Fisher information efficiency is 56\%; genome-wide aggregation across $\sim$300{,}000 windows yields $>$99.99\% information capture.

This framework explains the three key results. The 97.95\% simulated concordance confirms that assembly identity captures the essential TMRCA signal. The Platinum pedigree's 99.53\% mean concordance exploits the strong TMRCA contrast between unrelated grandparents, with only two meioses producing long, easily detected tracts. The remaining gap on real data (76.45\%) reflects a measurement limitation---per-window SNR of 0.69--1.06---not a theoretical one.

\subsection{Advantages of the Pangenome Approach}

Beyond the structural advantages outlined in Section~\ref{sec:introduction} (chromosome-scale phasing, all-variant capture, no ascertainment bias), a notable empirical finding is that IBD detection power scales with coalescent diversity ($\pi$) rather than genotyping platform coverage: SNR is highest in African populations (1.06) and lowest in East Asian (0.69), inverting the hierarchy of SNP array-based methods. This makes the approach particularly valuable for under-represented populations~\cite{lachance2012ascertainment}. The PAF-direct mode completes a full 22-autosome scan in 195\,s using 101\,MB RAM, making whole-genome analysis practical on a laptop.

\subsection{Algorithmic Maturity and the Emission Ceiling}

Our benchmarking program (241 benchmarks) reveals that the HMM operates near an empirical ceiling. The transition dominance ratio $R \approx 2{,}800$ places the system in a regime where emission improvements are bounded: a feature saturation theorem shows any emission transformation changes accuracy by $<$0.3\,pp, because only $\sim$6 boundary windows per chromosome are sensitive while $\sim$14{,}000 are locked by accumulated transition evidence. This explains why 32 independently motivated emission features do not compound. The pipeline retains $\sim$73\% of per-window Fisher information about TMRCA. All accuracy claims are grounded in empirical benchmarks; prescriptive predictions about intervention effects achieved 0/8 accuracy due to cascading interactions through auto-configuration (Supplementary Section~J.3).

\subsection{Limitations}

Several limitations should be noted. Per-window identity has a fundamentally low SNR for ancestry discrimination ($\sim$0.001 difference between populations at 10\,kb resolution), making the HMM dependent on temporal smoothing rather than per-window classification. Minority ancestry detection degrades as $\pi_j^{3/2}$ (Supplementary Section~J.8), explaining why AMR at 20\% has $\sim$1/4 the detection power of EUR at 50\%. The current HPRC panel (46 haplotypes) limits accuracy for under-represented populations, and panel scaling follows $O(1/\ln N)$---doubly logarithmic in panel size.

For IBD, the per-window identity ceiling limits segment-level concordance with hap-ibd (F1 = 0.514), though pair-level detection is strong (68.8\% recall, top pairs at Jaccard = 0.953). Identity-based IBD performs worse in high-SV regions due to alignment fragmentation, and logit-space emission transforms, while theoretically improving model adequacy, regress on real data due to bimodal emission artifacts (Supplementary Section~J.7).

Assembly quality and pangenome completeness remain prerequisites: regions without aligned segments produce gaps in the observation sequence, and assembly errors can mimic or mask IBD signal.

\subsection{Future Directions}

Several extensions follow from our findings. Graph-topology features (node sharing, bubble traversal patterns in pangenome graphs) may capture allele-specific information lost in aggregate identity, particularly in SV-rich regions where identity-based detection fails. Larger reference panels from HPRC Release~3 ($>$350 assemblies) will enable $K \geq 5$ population models and improve accuracy for under-represented groups. The per-window eGRM output enables downstream applications---local PCA, per-window $F_{ST}$, heritability estimation, and association testing---without VCF generation. Ancestry-specific eGRM~\cite{tang2025asegrm} conditioned on local ancestry calls provides an assembly-based as-eGRM without ARG inference. IBD segment lengths enable VCF-free demographic inference via established tools (HapNe-IBD, IBDNe), and a multi-pulse EM algorithm extending MultiWaver~\cite{ni2018multiwaver} is already implemented for ancestry tract-based dating. Finally, including archaic assemblies (Neanderthal, Denisovan) as reference states would enable assembly-based introgression detection.
